\documentclass[../DoAn.tex]{subfiles}
\begin{document}

\section{Next.js 14 -- Khung phát triển giao diện người dùng}
\label{section:3.1}

Next.js là khung phát triển ứng dụng web do Vercel phát hành trên nền React. Phiên bản 14 ra mắt tháng 11 năm 2023 đưa App Router lên thành mô hình định tuyến mặc định thay cho Pages Router của các phiên bản trước \cite{nextjsdocs}. App Router tổ chức ứng dụng dưới dạng cây thư mục trong \texttt{app/}: mỗi thư mục con ánh xạ thành một đoạn URL và có thể chứa các tệp \texttt{page.tsx}, \texttt{layout.tsx}, \texttt{loading.tsx}, \texttt{error.tsx} biểu diễn các trạng thái tương ứng.

Khác biệt kỹ thuật chính của App Router là Server Components: thành phần React mặc định chạy phía máy chủ, có thể truy vấn cơ sở dữ liệu hoặc gọi API nội bộ trực tiếp trong hàm render mà không cần state hay effect. Các thành phần cần xử lý sự kiện, hook hoặc API DOM khai báo \texttt{'use client'} ở đầu tệp. Sự phân tách này giảm khối lượng JavaScript gửi xuống trình duyệt và đơn giản hoá việc lấy dữ liệu cho trang chỉ đọc.

Đặc thù phân vùng theo bốn vai trò người dùng của PM QLKT khớp tự nhiên với App Router: các trang đặt lần lượt dưới \texttt{app/super-admin/}, \texttt{app/admin/}, \texttt{app/manager/}, \texttt{app/user/}; mỗi nhánh có \texttt{layout.tsx} riêng để áp các kiểm tra phân quyền chung. Trang Bảng điều khiển dành cho Admin tại \texttt{app/admin/dashboard/page.tsx} có thể gọi trực tiếp hàm tổng hợp số liệu phía máy chủ và trả về HTML đã render, tránh nhiều yêu cầu fetch nhỏ từ trình duyệt và rút ngắn thời gian hiển thị nội dung đầu tiên.

Đồ án đã cân nhắc hai phương án thay thế. Remix có mô hình data loader đơn giản hơn nhưng cộng đồng và hệ sinh thái plugin nhỏ hơn rõ rệt, đồng thời thiếu công cụ tách bundle theo route trưởng thành như Next.js. SvelteKit có hiệu năng runtime tốt nhưng buộc đội ngũ học thêm framework view mới ngoài React, trong khi Ant Design, thư viện UI chủ lực của hệ thống, mặc định phục vụ React. Vì vậy Next.js được chọn để tận dụng hệ sinh thái React sẵn có và mô hình Server Components mà không phải đánh đổi thư viện UI.

\section{TypeScript và Express -- Nền tảng phía máy chủ}
\label{section:3.2}

Express là khung HTTP tối giản trên nền Node.js, ra mắt từ năm 2010 và đến nay vẫn là một trong những dự án mã nguồn mở phổ biến nhất hệ sinh thái JavaScript \cite{expressdocs}. Tư tưởng chủ đạo là chuỗi middleware: mỗi yêu cầu đi qua một dãy hàm có chữ ký \texttt{(req, res, next)}, mỗi hàm hoặc thay đổi đối tượng yêu cầu hoặc kết thúc luồng. Mô hình này phù hợp với nhu cầu xác thực, ghi nhật ký và kiểm tra dữ liệu của PM QLKT.

TypeScript được dùng song song với Express ở toàn bộ backend để phát hiện lỗi từ thời điểm biên dịch. Cấu hình \texttt{tsconfig.json} đặt ở chế độ nới lỏng (\texttt{strict: false}, \texttt{strictNullChecks: false}) nhằm dung hoà giữa độ chặt của kiểu và sự tiện dụng khi thao tác với các đối tượng có nhiều trường tuỳ chọn của Express. Đổi lại, các kiểu Prisma sinh tự động từ schema được sử dụng xuyên suốt tầng Service và Repository, giữ tính nhất quán với cơ sở dữ liệu nhờ trình soạn thảo gợi ý đầy đủ.

Đặc tính chuỗi middleware được khai thác qua một mẫu năm tầng cố định cho route nhạy cảm: \texttt{verifyToken -> checkRole -> validate(schema) -> auditLog(options) -> controller.method}. Mỗi tầng đảm nhiệm một mối quan tâm độc lập như xác thực, phân quyền, kiểm tra dữ liệu và ghi nhật ký, cho phép tái sử dụng giữa các route. Hàm \texttt{catchAsync} gói handler bất đồng bộ để chuyển lỗi sang middleware xử lý lỗi tập trung, tránh viết lặp \texttt{try/catch} ở mỗi controller.

Hai phương án backend khác cũng được cân nhắc. Fastify nhanh hơn Express khoảng 20\,\% qua benchmark độc lập nhưng dùng schema-based validation gắn cứng với JSON Schema, gây trùng lặp với Zod đã chọn cho frontend (Mục~\ref{section:3.7}). NestJS cung cấp kiến trúc module hoá rõ ràng với decorator nhưng buộc đội ngũ học một framework áp đặt nhiều quy ước và tăng số lượng lớp trừu tượng như Module, Provider, Guard. Chi phí học này không tương xứng với quy mô một đồ án 23 model. Express được chọn vì middleware chain đủ linh hoạt cho mẫu năm tầng nói trên, hệ sinh thái plugin dày (Helmet, express-rate-limit, multer đều được dùng), và dễ refactor sang Fastify hoặc NestJS sau này nếu nhu cầu thay đổi.

\section{PostgreSQL và Prisma ORM}
\label{section:3.3}

PostgreSQL là hệ quản trị cơ sở dữ liệu quan hệ mã nguồn mở phát triển từ năm 1986, hiện được duy trì bởi cộng đồng phát triển PostgreSQL \cite{pgdocs}. Hệ này được chọn nhờ độ ổn định cao, tuân thủ chuẩn SQL chặt chẽ và hỗ trợ nhiều kiểu dữ liệu nâng cao như JSON, JSONB, mảng, kiểu liệt kê và kiểu thời gian có múi giờ, tất cả đều được PM QLKT sử dụng.

Prisma là ORM thế hệ mới, không sinh mã kiểu Active Record mà sinh ra một Client typed an toàn từ tệp khai báo \texttt{schema.prisma} duy nhất \cite{prismadocs}. Mỗi thay đổi schema được áp dụng qua \texttt{prisma migrate dev}, sinh các tệp SQL trong \texttt{prisma/migrations/} và đẩy lên kho mã nguồn để theo vết lịch sử. Mọi truy vấn \texttt{findUnique}, \texttt{findMany}, \texttt{create} đều trả về kiểu suy luận trực tiếp từ schema, có gợi ý đầy đủ tại trình soạn thảo.

Schema của PM QLKT gồm 23 model trong 618 dòng. Quy ước đặt tên giữ nguyên tên thực thể tiếng Việt ở cả tầng mã lẫn tầng cơ sở dữ liệu: tên bảng vật lý và tên model dùng \texttt{PascalCase} tiếng Việt (\texttt{QuanNhan}, \texttt{DanhHieuHangNam}, \texttt{HoSoCongHien}, \texttt{BangDeXuat}); tên cột và field dùng \texttt{snake\_case} tiếng Việt (\texttt{ho\_ten}, \texttt{ngay\_sinh}, \texttt{so\_quyet\_dinh}, \texttt{quan\_nhan\_id}). Khi đọc lệnh \texttt{prisma.quanNhan.findMany(\dots)}, lập trình viên xác định ngay bảng được truy vấn mà không phải tra cứu mapping. Cột giữ \texttt{snake\_case} đồng nhất với truyền thống đặt tên dữ liệu hành chính tiếng Việt và tương thích với các báo cáo SQL viết tay. Bảng \ref{tab:so-sanh-orm} so sánh ngắn gọn ba ORM phổ biến để làm rõ lý do chọn Prisma.

\begin{table}[H]
\centering
\caption{So sánh ba ORM phổ biến trên Node.js}
\label{tab:so-sanh-orm}
\begin{tabular}{|p{4cm}|p{3cm}|p{3cm}|p{4cm}|}
\hline
\textbf{Tiêu chí} & \textbf{Sequelize} & \textbf{TypeORM} & \textbf{Prisma} \\ \hline
Năm ra mắt & 2010 & 2016 & 2019 \\ \hline
Kiểu định nghĩa schema & JavaScript / mô hình lớp & Decorator + lớp & DSL trong tệp \texttt{schema.prisma} \\ \hline
Hỗ trợ TypeScript & Có nhưng cần khai báo bổ sung & Tốt & Sinh kiểu hoàn toàn tự động \\ \hline
Migration & Có (kém trực quan) & Có & Sinh tự động và có thể đảo ngược \\ \hline
Hiệu suất truy vấn & Trung bình & Trung bình & Tối ưu nhờ engine Rust \\ \hline
Cộng đồng & Lớn nhưng đang chậm phát triển & Trung bình & Đang tăng nhanh \\ \hline
Phù hợp cho dự án có nhiều quan hệ & Cần cấu hình thủ công & Cần decorator nhiều & Khai báo gọn, sinh truy vấn join hiệu quả \\ \hline
\end{tabular}
\end{table}

Với 23 model và nhiều quan hệ phức tạp (vd: \texttt{BangDeXuat <-> QuanNhan} qua trường \texttt{data\_danh\_hieu} JSONB hoặc qua \texttt{HoSoNienHan}), khả năng sinh truy vấn \texttt{include}, \texttt{select} của Prisma giảm rõ rệt khối lượng mã ở tầng Repository so với SQL viết tay. Đây là lý do trọng yếu để PM QLKT chọn Prisma.

\section{Ant Design và Tailwind CSS -- Bộ giao diện}
\label{section:3.4}

Frontend của PM QLKT kết hợp hai thư viện giao diện ở hai vai trò khác nhau. Ant Design là thư viện thành phần React do Alibaba phát hành \cite{antdesigndocs}, cung cấp nhiều thành phần phức tạp đã hoàn thiện như \texttt{Table} có lọc và phân trang, \texttt{Form} tích hợp xác thực, \texttt{Select} đa lớp, \texttt{DatePicker} hỗ trợ múi giờ, \texttt{Modal} lồng nhau và \texttt{Notification}. Đây là thư viện chủ lực của hệ thống: các biểu mẫu đề xuất nhiều bước, bảng danh sách quân nhân có cột động và các hộp thoại xác nhận phê duyệt đều dùng Ant Design.

Tailwind CSS là khung CSS utility-first \cite{tailwinddocs} cho phép áp các lớp đơn lẻ trực tiếp trong JSX để mô tả bố cục, khoảng cách, màu sắc và đáp ứng nhiều kích cỡ màn hình. Tailwind không thay thế Ant Design mà bổ sung tại những điểm cần tinh chỉnh bố cục ngoài phạm vi của các thành phần dựng sẵn, ví dụ căn lề giữa các thẻ thống kê trên Bảng điều khiển, hay xếp lưới hai cột trên màn hình lớn và một cột trên thiết bị di động.

Để quản lý các lớp Tailwind theo điều kiện, hệ thống dùng tiện ích \texttt{cn()} kết hợp \texttt{clsx} và \texttt{tailwind-merge}, giúp gộp và khử trùng lặp lớp CSS khi dựng thành phần. Quy ước nội bộ là Ant Design cho biểu mẫu và bảng nghiệp vụ, còn Tailwind cho bố cục và tinh chỉnh giao diện ngoài phạm vi của các thành phần dựng sẵn.

\section{Socket.IO -- Kênh thông báo thời gian thực}
\label{section:3.5}

Socket.IO trừu tượng hoá giao thức WebSocket và tự động chuyển sang HTTP long polling khi mạng giữa máy khách và máy chủ không cho phép kết nối WebSocket trực tiếp \cite{socketiodocs}. Thư viện cung cấp khái niệm ``phòng'' (room) cho phép phát thông điệp tới một nhóm kết nối được gán nhãn, phù hợp với mô hình thông báo theo người dùng hoặc theo đơn vị.

PM QLKT dùng Socket.IO cho hai trường hợp. Một là, khi Admin nhập tệp Excel hàng trăm bản ghi, máy chủ phát sự kiện tiến độ theo phần trăm hoàn thành để giao diện cập nhật thanh tiến trình, tránh polling từ phía trình duyệt. Hai là, khi đề xuất được phê duyệt, máy chủ phát thông điệp tới các kết nối Manager của đơn vị liên quan; giao diện hiển thị thông báo dạng pop-up (toast notification) ở góc màn hình mà không phải làm mới trang. Kết nối Socket.IO xác thực bằng JWT lấy từ \texttt{localStorage}, truyền vào tuỳ chọn \texttt{auth: \{ token \}} ở giai đoạn handshake; máy chủ dùng cùng \texttt{JWT\_SECRET} với REST để xác minh, giữ logic xác thực hợp nhất giữa REST và WebSocket.

\section{JSON Web Token với cơ chế làm mới luân phiên}
\label{section:3.6}

Hệ thống áp dụng JWT theo chuẩn RFC 7519 \cite{rfc7519} cho xác thực không trạng thái. Mỗi phiên đăng nhập sinh hai token: Access Token thời hạn ngắn (30 phút) ký bằng \texttt{JWT\_SECRET}, chứa định danh người dùng, mã vai trò và liên kết quân nhân, gửi kèm mỗi yêu cầu HTTP qua header \texttt{Authorization: Bearer}; Refresh Token thời hạn dài (7 ngày) ký bằng \texttt{JWT\_REFRESH\_SECRET} riêng, chỉ dùng để cấp Access Token mới khi token hiện tại hết hạn. Frontend lưu cả hai token vào \texttt{localStorage}; backend lưu thêm Refresh Token vào cột \texttt{TaiKhoan.refreshToken} để giữ khả năng thu hồi từ phía máy chủ.

Cơ chế làm mới luân phiên (refresh token rotation) phát hành đồng thời cặp Access + Refresh Token mới ở mỗi lần làm mới và ghi đè cột \texttt{refreshToken}; lần dùng Refresh Token cũ sau đó đều bị từ chối vì không còn khớp với giá trị trong cơ sở dữ liệu. Khi Refresh Token bị đánh cắp, chỉ bên dùng token đầu tiên thực hiện được lần làm mới, bên dùng token sau đó sẽ thất bại dù là kẻ tấn công hay người dùng hợp pháp. Đồng thời quản trị viên có thể buộc đăng xuất bằng cách xoá cột \texttt{refreshToken}. Hai khoá bí mật riêng biệt cho hai loại token theo nguyên tắc phòng thủ nhiều lớp: nếu một khoá bị lộ, chỉ một loại token bị ảnh hưởng. Mật khẩu trước khi lưu vào \texttt{TaiKhoan} được băm bằng bcrypt với hệ số chi phí 10 \cite{bcrypt}. Hệ số này cân bằng giữa thời gian xác thực, khoảng 80\,ms trên máy chủ phát triển, và độ khó vét cạn nếu cơ sở dữ liệu bị rò rỉ.

\section{Zod -- Kiểm tra hợp lệ dữ liệu hai phía}
\label{section:3.7}

PM QLKT thống nhất dùng Zod \cite{zoddocs} làm thư viện xác thực dữ liệu cho cả backend lẫn frontend, thay vì tách Joi cho backend và Zod cho frontend như nhiều dự án Express truyền thống. Tại backend, mọi route nhận dữ liệu từ client đều đi qua middleware \texttt{validate(schema)} dùng \texttt{ZodType} để kiểm tra \texttt{req.body}, \texttt{req.query} và \texttt{req.params}; schema tập trung tại \texttt{src/validations/}, mỗi miền nghiệp vụ một tệp \texttt{*.validation.ts}. Phương thức \texttt{.strict()} của Zod từ chối những trường không khai báo trong schema, ngăn người dùng đẩy lên các thuộc tính ngoài ý muốn.

Tại frontend, các biểu mẫu Ant Design Form được xác thực bằng Zod kết hợp React Hook Form. Zod tích hợp chặt với TypeScript: kiểu của giá trị form suy luận trực tiếp từ schema qua \texttt{type FormValues = z.infer<typeof formSchema>}, tránh duy trì song song interface và schema. Việc dùng chung một thư viện ở hai phía giảm chi phí học và cho phép viết các đoạn ràng buộc tương tự (vd: CCCD 12 chữ số, năm hợp lệ) bằng cùng một cú pháp, mở đường cho việc trích xuất schema dùng chung sang gói chia sẻ trong tương lai. Logic xác thực không tin tưởng tầng frontend: backend luôn kiểm tra lại dữ liệu trước khi ghi vào cơ sở dữ liệu.

\section{ExcelJS -- Công cụ nhập và xuất Excel}
\label{section:3.8}

ExcelJS \cite{exceljsdocs} là thư viện JavaScript cho phép đọc, ghi và thao tác workbook \texttt{.xlsx} phía máy chủ Node.js. Khác với các thư viện chỉ sinh CSV hay HTML giả lập Excel, ExcelJS hỗ trợ định dạng ô (font, màu nền, viền), khoá ô để chặn chỉnh sửa cột hệ thống, công thức, và quan trọng nhất là Data Validation cho phép tạo danh sách thả xuống ngay trong tệp mẫu xuất ra.

ExcelJS phục vụ hai chiều dữ liệu trong PM QLKT. Khi xuất tệp mẫu nhập liệu cho mỗi nhóm khen thưởng, máy chủ sinh workbook với các sheet \texttt{QuanNhan}, \texttt{DanhHieuHangNam}, \texttt{ThanhTichKhoaHoc}, \dots; mỗi sheet cố định cấu trúc cột, áp Data Validation cho cột danh hiệu để chỉ chấp nhận các giá trị hợp lệ và đính chú thích tiếng Việt giải thích từng cột. Khi nhập tệp Excel do người dùng tải lên, máy chủ đọc tuần tự từng dòng, đối chiếu với schema Zod và các ràng buộc Prisma, rồi gom kết quả vào màn hình xem trước trước khi ghi xuống cơ sở dữ liệu. Bước xác nhận chạy trong một transaction Prisma để đảm bảo nguyên tắc all-or-nothing: nếu có bất kỳ dòng nào gặp lỗi tại thời điểm ghi, toàn bộ thao tác được hoàn tác.

\end{document}
